\section{Introduction}
Pfaffian functions satisfy triangular systems of first-order partial differential equations with polynomial coefficients. They extend polynomial functions to include transcendental examples such as $\exp$, $\log$ on its domain, and $\tanh$, while retaining effective quantitative bounds on the geometry of their zero sets.

This paper has three main aims. First, we establish local tube-volume bounds for regular Pfaffian hypersurfaces, including unbounded ones. Second, we apply these bounds to the robustness of neural network classifiers, using a condition number defined by distance to the decision boundary. Third, for a single sigmoid layer with rational first-layer weights, we replace the exponential dependence on width in the general Pfaffian estimate by a polynomial bound.

Let $f\colon \R^n\to \R$ be a Pfaffian function with chain length $s$, chain degree $\alpha$ and degree $\beta$ (as defined in Section~\ref{sec:pfaffian}), where $\alpha,\beta\ge1$, and suppose $0$ is a regular value of $f$. For $V=\mathcal Z(f)$ and $X$ uniform in any ball $B(p,\rho)$, Theorem~\ref{thm:prob_bound} gives
\[
\prob\{\dist(X,V)\le\varepsilon\}
\le C_{\alpha,\beta,s,n}
\left[\left(1+(\alpha+\beta)\frac{\varepsilon}{\rho}\right)^n
-\left(1+\frac{\varepsilon}{\rho}\right)^n\right],
\]
where the explicit constant~\eqref{eq:pfaffian-tube-constant} depends only on the format and dimension. For fixed parameters the right-hand side is $O(\varepsilon/\rho)$ as $\varepsilon/\rho\to0$. The proof combines the integral-geometric tube estimate of~\cite{lotz2015volume} with Khovanskii's zero bound~\cite{khovanskii1991fewnomials}. 
%Truncating just beyond the region in which nearest points can lie removes the need for a separate boundary contribution.

Many smooth neural network activation functions, including the logistic sigmoid $\sigma(t)=(1+\econst^{-t})^{-1}$ and $\tanh$, are Pfaffian. A neural network classifier $F\colon\R^n\to\R^m$ assigns an input $x\in \R^n$ to a label, or class, $k$, if $k\in \argmax_j F_j(x)$ (with ties broken arbitrarily), where each $F_j$ is referred to as a score function. Networks built from Pfaffian activations and Pfaffian output maps are Pfaffian, and the decision boundaries for classification problems are semi-Pfaffian. Let $\Sigma$ be the locus of ties for the maximal score. We define the local condition number
\[
\mC_p(x)=\frac{\|x-p\|}{\dist(x,\Sigma)}.
\]
When all pairwise score differences have $0$ as a regular value, the denominator is the infimum of perturbation sizes that change a uniquely predicted label. This measures label stability, independently of whether the label agrees with ground truth. The tube estimate gives explicit $O(1/t)$ bounds for $\prob\{\mC_p(X)>t\}$ both for uniform sampling in $B(p,\rho)$ and for Gaussian sampling $X\sim\mathcal N(p,\varsigma^2 I_n)$ at every positive scale (Theorems~\ref{thm:condition-uniform} and~\ref{thm:condition-gaussian}).

For a sigmoid network with $h$ hidden units, the general bound contains the Khovanskii factor $2^{h(h-1)/2}$. For a single hidden layer,
we improve this substantially. Suppose its width is $w$, its first-layer weight vectors are $a_k\in\Q^n$, and its outputs are affine, optionally followed by softmax. Let $q\in\N_{>0}$ clear the denominators and put $L=\max_k\|qa_k\|_\infty\ge1$. If every pairwise logit difference has $0$ as a regular value, Corollary~\ref{cor:nn-condition-single-layer} gives
\[
\prob\{\mC_p(X)>t\}
\le\min\left\{1,\binom m2 H(n,L)\frac{w^n}{t}\right\},
\qquad t>0,
\]
for an explicit $H(n,L)$, for both sampling models and all centres and positive scales. The polynomial width dependence is at fixed $n,L$. 

Single-hidden-layer networks are a natural setting in which to study width. The universal approximation theorems of Cybenko~\cite{cybenko1989approximation} and Hornik, Stinchcombe, and White~\cite{hornik1989multilayer} show that finite sums of sigmoids are dense in continuous functions on compact sets. Barron's theorem~\cite{barron1993universal} gives an $L^2$-norm approximation error of order $C_f/\sqrt w$ for target functions with finite first Fourier moment $C_f$; see also~\cite{pinkus1999approximation}.

The geometric estimates use two representations of a scalar network
\[
f(x)=c_0+\sum_{k=1}^w d_k\sigma(a_k^{\trans}x+b_k).
\]
The tubular neighbourhood bounds depend on the maximal degree $\md(\mathcal Z(f))$ of the generalised Gauss map and on the section degrees $\md_i(\mathcal Z(f))$, defined in Section~\ref{sec:tubular}.
After the substitution $y_j=\econst^{-x_j/q}$, clearing a single product of positive denominators turns its Gauss-fibre equations determining the maximal degree into Laurent polynomial equations with a common zonotope containing their supports. Bernstein's theorem gives
\[
\md(\mathcal Z(f))\le n!L^nw^n.
\]
A compact grid construction shows that the exponent $n$ in this Gauss-map degree bound is optimal (Proposition~\ref{prop:single-layer-lowerbound}). To handle arbitrary affine sections, we keep the $n$ coordinate exponentials as a global Pfaffian chain and pull them back to the section. Its length is independent of $w$, so the standard Khovanskii theorem gives polynomial bounds for all section degrees. Finally, a separate zero count along lines bounds the leading tube coefficient by $O_{n,L}(w^n)$. Extending the degree estimate to greater depth is the subject of Conjecture~\ref{conj:multi-layer-degree}.

A limit of the regular levels $f=\pm\delta$ also gives a full-tube probability bound of $O_{n,L}(w^n\varepsilon/\rho)$ for possibly singular shallow zero sets (Corollary~\ref{cor:single-layer-singular-linear}).

\subsection{Previous and related work}
The problem of bounding the volume of a tubular neighbourhood of a set in terms of its geometric complexity has a long and rich history. In 1840, Steiner~\cite{steiner1840bestimmung} observed that the volume of the $\varepsilon$-neighbourhood of a convex body in $\R^3$ is a cubic polynomial in $\varepsilon$. A celebrated generalisation is due to Weyl~\cite{weyl1939volume}, who showed that the volume of the $\varepsilon$-tubular neighbourhood ($\varepsilon$ small enough) of a compact Riemannian submanifold of $\R^n$ is a polynomial in $\varepsilon$ whose coefficients are intrinsic curvature invariants; see also Hotelling~\cite{hotelling1939tubes} and the monograph by Gray~\cite{gray2004tubes}. Tube formulae entered numerical analysis through the work of Smale, Kostlan, Renegar, and Demmel, among others, who were interested in the probabilistic analysis of condition numbers (see the references in~\cite{burgisser2013condition,lotz2015volume,basu2021hausdorff}). The key observation is that if the set of ill-posed inputs of a numerical problem can be described as a subset of an algebraic variety, then a bound on the volume of its tubular neighbourhood directly translates into a bound on the probability distribution of the condition number.
Bounds on tubular neighbourhoods have been extended to singular algebraic sets by Basu and Lerario~\cite{basu2021hausdorff}.
In a complementary direction, Zhang and Kileel~\cite{zhang2025covering} bound the covering numbers of real algebraic varieties, images of polynomial maps, and general semialgebraic sets via a slicing argument that counts the connected components of generic affine sections, and deduce volume bounds for tubular neighbourhoods of such sets without any smoothness assumptions. Their bounds capture the correct leading order in $\varepsilon$, but are obtained from coverings by balls and therefore carry dimension-exponential constants rather than the degree-graded coefficients of the Weyl-type tube formula used here; the component counts of generic affine sections that control their estimates play a role analogous to the section degrees of the Gauss map that govern our tube formula, and admit Khovanskii-type bounds in the Pfaffian setting (Section~\ref{sec:pfaffian}). Among their applications are generalization bounds for deep networks with rational or ReLU activations, which measure the complexity of a network class in function space and are thereby complementary to our robustness analysis, which concerns the geometry of the decision boundary in input space.

Neural networks have been studied in the context of Pfaffian functions since the work of Macintyre and Sontag~\cite{macintyre1993finiteness}, who established finiteness results for sigmoidal networks. Karpinski and Macintyre~\cite{karpinski1997polynomial} showed that the VC dimension of sigmoidal Pfaffian networks is polynomial in the number of parameters; this was recently extended by D'Inverno, Bianchini, and Scarselli~\cite{dinverno2024vc} to graph neural networks with general Pfaffian activations. Bianchini and Scarselli~\cite{bianchini2014complexity} bound the sums of Betti numbers of decision regions for networks with Pfaffian activation functions. In particular, their Theorem~5 constructs, for every $\ell\ge1$, a sigmoid network with $\ell$ hidden layers and two neurons per hidden layer whose nonnegative decision region $\{x\colon f(x)\ge0\}$ has at least $2^{\ell-1}+1$ connected components. Their proof in Appendix~F iterates a fixed scalar map realised by one hidden sigmoid layer, establishing exponential growth of decision-region topology with depth even at width two.

The robustness of neural network classifiers with respect to adversarial perturbations has been of practical concern in contexts ranging from spam filtering to medical diagnosis, and a systematic investigation was initiated by Szegedy et al.~\cite{szegedy2013intriguing}. Subsequent work developed efficient attacks and robustness surrogates based on first-order or local geometric approximations, including the fast-gradient method of Goodfellow, Shlens, and Szegedy~\cite{goodfellow2015explaining}, the DeepFool algorithm of Moosavi-Dezfooli, Fawzi, and Frossard~\cite{moosavi2016deepfool}, and the analysis of adversarial, random, and semi-random perturbations by Fawzi, Moosavi-Dezfooli, and Frossard~\cite{fawzi2016robustness}. Another theoretical line explains adversarial vulnerability through high-dimensional geometry and concentration of measure; see, for example, Fawzi, Fawzi, and Fawzi~\cite{fawzi2018adversarial} and Mahloujifar, Diochnos, and Mahmoody~\cite{mahloujifar2019curse}. The robustness of classifiers is an inherently geometric problem. A classifier $F\colon \R^n\to \R^m$ partitions the input space into decision regions separated by a decision boundary $\Sigma$. Under the pairwise regularity assumptions used here, the distance $\dist(x,\Sigma)$ is the infimum of perturbation sizes that change the predicted label of an input~$x$ with a unique maximal score, and one can define a condition number $\mC(x) = \|x\|/\dist(x,\Sigma)$ by analogy with the theory of conditioning in numerical analysis~\cite{burgisser2013condition}. Bounding the volume of the tubular neighbourhood of $\Sigma$ therefore provides quantitative robustness guarantees and, more precisely, tail bounds on the probability distribution of the condition number.

In the case of neural networks with piecewise-linear activations such as ReLU, the decision boundary is itself piecewise-linear, with combinatorial complexity controlled by the linear-region structure induced across the layers. Mont\'ufar, Pascanu, Cho, and Bengio~\cite{montufar2014} initiated the systematic counting of linear regions, proving in particular exponential lower bounds for deep networks; upper bounds of order $O_{n,\ell}\!\bigl(\prod_j n_j^n\bigr)$ for an $\ell$-layer ReLU network of widths $n_1, \ldots, n_\ell$ at fixed input dimension and depth were obtained by Raghu, Poole, Kleinberg, Ganguli, and Sohl-Dickstein~\cite{raghu2017} and Serra, Tjandraatmadja, and Ramalingam~\cite{serra2018}, with refined counts by Hanin and Rolnick~\cite{hanin2019}. A different perspective due to Zhang, Naitzat, and Lim~\cite{zhang2018tropical} interprets ReLU networks as tropical rational maps and relates their linear-region structure to polyhedral subdivisions associated with a tropical-polynomial representation.

A recent program of neuroalgebraic geometry~\cite{marchetti2025neuroalgebraic} studies the function space of a network, the \emph{neuromanifold} swept out as the weights vary, as a (semi-)algebraic variety, relating its algebraic invariants such as dimension, degree, and singularities to learning-theoretic properties. This shares our outlook of analysing neural networks through real-geometric invariants, including degree notions, but the object and regime differ: that program concerns the function space of polynomial networks, where algebraic geometry applies directly, whereas we study the decision boundary in input space of networks with transcendental Pfaffian activations such as the sigmoid, for which the relevant tame structure is o-minimal and the natural finiteness tools are Khovanskii's fewnomial theory and the Bernstein--Kushnirenko--Khovanskii theorem. 

\subsection{Structure of the paper}
Section~\ref{sec:pfaffian} reviews the necessary background on Pfaffian functions and Pfaffian sets in a self-contained manner. Section~\ref{sec:tubular} derives the tube formula for smooth Pfaffian hypersurfaces and applies it to obtain bounds on the volume of their tubular neighbourhoods. In Section~\ref{sec:tubular-neural}, we apply the tube formula to obtain bounds on the volume of tubular neighbourhoods of the decision boundary of neural network classifiers with Pfaffian activations. We also present our second main result, the tube formula for one-layer sigmoid networks that is polynomial in the width of the network. Section~\ref{sec:robustness} applies the tube formula to obtain tail bounds on the probability distribution of the condition number of neural network classifiers with Pfaffian activations. Section~\ref{sec:conclusions} outlines some future directions.

\subsection{Acknowledgements} The authors are grateful to Abhiram Natarajan for many insightful discussions on Pfaffian geometry and its applications. P.L. was funded by the London Mathematical Society through the Undergraduate Research Bursary URB-2022-69, and by a London School of Geometry and Number Theory--Imperial College London/King's College London/University College London PhD studentship, which is supported by the Engineering and Physical Sciences Research Council [EP/S021590/1]. M.L. was funded by EPSRC Grant EP/W00383X/1. 

\subsection{Notation and conventions}\label{subsec:notation}
We denote by $\N$ the natural numbers including $0$, by $\N_{>0}$ the positive integers, and by $[n]=\{1,\dots,n\}$, with $[0]=\varnothing$. The space $\R[X_1,\dots,X_n]_{\le d}$ consists of polynomials of total degree at most $d$. We identify a polynomial with its associated function when the domain is clear, and write $\mathcal Z(f_1,\dots,f_k)$ for their common zero set (also for nonpolynomial functions).

We use the Euclidean inner product $\langle\cdot,\cdot\rangle$ and norm $\|\cdot\|$, and write $\|x\|_\infty=\max_j|x_j|$. The vectors $e_j$ form the standard basis, $A^{\trans}$ denotes transpose, and $I_d$ is the $d\times d$ identity matrix. We write $B(p,\rho)$ for the closed Euclidean ball, $S^{n-1}(p,\rho)=\partial B(p,\rho)$, $S^{n-1}=S^{n-1}(0,1)$, and $\omega_n=\vol B(0,1)$, where $\vol$ is Euclidean volume in the ambient space. Also, $\dist(x,A)=\inf_{a\in A}\|x-a\|$, with $\dist(x,\varnothing)=\infty$.

For a smooth map $F$, $\diff{F}(x)$ denotes its differential at $x$. For a smooth scalar function $f$ and a constant vector $\tau$ in its domain space, we write
\[
\partial_\tau f(x):=\left.\frac{\diff{}}{\diff t}f(x+t\tau)\right|_{t=0}
=\diff{f}(x)[\tau]=\langle\nabla f(x),\tau\rangle.
\]
The vector $\tau$ need not be a unit vector. In coordinates, $\partial_j f=\partial_{x_j}f=\partial_{e_j}f$, $\partial_{jk}f=\partial_j\partial_k f$, and $\operatorname{Hess}f=(\partial_{jk}f)_{j,k}$. A numerical subscript refers to the corresponding coordinate of the function's current argument. A variable subscript, as in $\partial_{y_j}$ or $\partial_{t_k}$, specifies the coordinate explicitly.

We reserve $\sigma$ for the logistic sigmoid (and $\sigma^i_j$ for general hidden activations), and use $\varsigma>0$ for a Gaussian sampling scale: $\mathcal N(p,\varsigma^2 I_n)$. Subscripts on $O$ and $\Omega$ list the parameters on which their implicit constants may depend.
